\section{Framework Description} \label{sec:framework}

This section describes the architecture, components, and algorithms of the Red-Team vs.\ Blue-Team Evaluation Framework.
The framework operates as a closed-loop system: a simulated autonomous platform generates normal MAVLink telemetry, red-team attack modules inject adversarial traffic, a blue-team detector analyses the combined stream, and a centralised orchestrator coordinates timing, logging, and metric collection.

\subsection{Problem Definition}

We define the red-team vs.\ blue-team evaluation problem as follows.

Let $\mathcal{P}$ be a \emph{platform} that produces a stream of MAVLink messages $m_1, m_2, \dots$ over time.
Each message $m_t$ has a type $\tau(m_t)$, a source identifier $\sigma(m_t)$, and a timestamp $t$.

An \emph{attack catalogue} $\mathcal{A} = \{a_1, \dots, a_k\}$ is a set of attack modules.
Each attack $a_i$ is parameterised by a target address, a message rate $r_i$, and a duration $d_i$.
When active, $a_i$ injects additional messages into the same channel that $\mathcal{P}$ uses.

A \emph{schedule} $\mathcal{S}$ assigns attacks to time intervals.
In static mode, attacks run sequentially at fixed rates.
In randomised mode, $\mathcal{S}$ is sampled from a distribution controlled by a seed $s$: the attack order is shuffled, each episode is split into segments with independently sampled rates, and random pauses separate consecutive episodes.

A \emph{detector} $\mathcal{F}$ maps each message to an anomaly score $\mathcal{F}(m_t) \in [0, 1]$, where 1 is normal and 0 is maximally anomalous.
The detector first trains on a baseline window of benign traffic from $\mathcal{P}$ and then scores all subsequent messages.

A \emph{defence} $\mathcal{D}$ receives the anomaly score (or a related signal such as the Markov transition probability) and produces a binary decision $\mathcal{D}(m_t) \in \{\textit{block}, \textit{pass}\}$.

The \emph{orchestrator} $\mathcal{O}$ coordinates $(\mathcal{P}, \mathcal{A}, \mathcal{S}, \mathcal{F}, \mathcal{D})$.
It launches $\mathcal{P}$, starts $\mathcal{F}$ in a background thread, waits for baseline training to complete, executes $\mathcal{S}$, and after all attacks finish it computes the metric vector:
\begin{equation}
\mathcal{M} = \bigl(\textit{precision},\; \textit{recall},\; F_1,\; t_{\textit{det}},\; \beta,\; \mu\bigr)
\end{equation}
where $t_{\textit{det}}$ is the detection latency (time from the first attack to the first alert), $\beta$ is the block rate (fraction of attack messages stopped by $\mathcal{D}$), and $\mu$ is the mission impact ratio (fraction of monitoring time during which the system operated in a degraded state).

The evaluation goal is: given a fixed $\mathcal{P}$ and $\mathcal{A}$, compare different combinations of $\mathcal{F}$, $\mathcal{D}$, and $\mathcal{S}$ to identify which defence strategy minimises $\mu$ while maintaining high precision and recall.
Because $\mathcal{O}$ controls all variables and produces $\mathcal{M}$ in a standardised format, results across different configurations are directly comparable without external coordination.

\subsection{System Architecture}

The framework comprises four main components connected via UDP on port 14550 (the standard MAVLink ground-station port):

\begin{enumerate}
    \item \textbf{Platform Simulator} (\texttt{sitl\_sim.py}): Emits normal MAVLink messages (HEARTBEAT, PING, PARAM\_REQUEST\_LIST, STATUSTEXT) at a configurable rate (default 10~msg/s) from source system ID~1, simulating a quadrotor in active flight.

    \item \textbf{Red-Team Attack Catalogue} (\texttt{attacks/}): Six independent attack modules, each configurable in rate (msg/s) and duration (seconds), injecting malicious MAVLink packets from rogue source system IDs.

    \item \textbf{Blue-Team Detector and Defence} (\texttt{detector/}): A two-stage pipeline consisting of (a)~an 11-feature extraction engine, (b)~an Isolation Forest anomaly scorer, (c)~a Markov transition model, and (d)~an optional active defence that blocks suspicious messages.

    \item \textbf{Orchestrator} (\texttt{run\_demo.py}): The central coordinator that launches the simulator, starts the detector, schedules attacks according to a deterministic or randomised plan, collects results, and produces standardised metrics and plots.
\end{enumerate}

Figure~\ref{fig:framework} shows the data flow between these components.
The orchestrator sits above all other modules and controls their lifecycle, but it never accesses their internal state.
Normal and malicious traffic merge on the same UDP channel before reaching the blue-team pipeline, which processes every message identically regardless of its origin.

\begin{figure*}[t]
\centering
\resizebox{\textwidth}{!}{%
\begin{tikzpicture}[
    node distance=1.2cm and 1.8cm,
    box/.style={draw, rounded corners, minimum width=2.4cm, minimum height=0.9cm, align=center, font=\small},
    bigbox/.style={draw, rounded corners, minimum width=3.2cm, minimum height=0.9cm, align=center, font=\small},
    orch/.style={draw, rounded corners, fill=gray!15, minimum width=4cm, minimum height=1cm, align=center, font=\small\bfseries},
    redbox/.style={draw, rounded corners, fill=red!10, minimum width=2.2cm, minimum height=0.8cm, align=center, font=\scriptsize},
    bluebox/.style={draw, rounded corners, fill=blue!10, minimum width=2.4cm, minimum height=0.8cm, align=center, font=\scriptsize},
    simbox/.style={draw, rounded corners, fill=green!10, minimum width=3cm, minimum height=0.9cm, align=center, font=\small},
    outbox/.style={draw, rounded corners, fill=yellow!12, minimum width=2.2cm, minimum height=0.8cm, align=center, font=\scriptsize},
    arr/.style={-{Stealth[length=2.5mm]}, thick},
    darr/.style={-{Stealth[length=2.5mm]}, thick, dashed},
    label/.style={font=\scriptsize\itshape, text=gray!70!black},
]

% ---- Orchestrator (top) ----
\node[orch] (orch) {Orchestrator\\[-1pt]\texttt{\scriptsize run\_demo.py}};

% ---- Simulator (left) ----
\node[simbox, below left=1.6cm and 3.5cm of orch] (sim) {Platform Simulator\\[-1pt]\texttt{\scriptsize sitl\_sim.py}};

% ---- Red Team (left column) ----
\node[redbox, below=2.8cm of sim, xshift=-1.2cm] (atk1) {Heartbeat\\Flood};
\node[redbox, right=0.3cm of atk1] (atk2) {Ping\\Flood};
\node[redbox, right=0.3cm of atk2] (atk3) {Param-Req\\Flood};
\node[redbox, below=0.3cm of atk1] (atk4) {MITM\\Spoof};
\node[redbox, right=0.3cm of atk4] (atk5) {Replay\\Pattern};
\node[redbox, right=0.3cm of atk5] (atk6) {Command\\Injection};

% Red team container
\begin{scope}[on background layer]
\node[draw=red!60, fill=red!3, rounded corners, fit=(atk1)(atk2)(atk3)(atk4)(atk5)(atk6),
      inner sep=0.35cm, label={[font=\small\bfseries, text=red!70!black]above:Red Team (Attacks)}] (redgroup) {};
\end{scope}

% ---- UDP Channel (center) ----
\node[draw, fill=orange!8, rounded corners, minimum width=1.6cm, minimum height=3.2cm,
      right=2.2cm of redgroup, align=center, font=\small] (udp) {UDP\\Port\\14550\\[-2pt]\scriptsize(MAVLink)};

% ---- Blue Team (right) ----
\node[bluebox, right=2.2cm of udp, yshift=1.2cm, minimum width=2.8cm] (feat) {Feature Engine\\[-1pt]\scriptsize 11-dim vector};
\node[bluebox, below=0.4cm of feat, minimum width=2.8cm] (iforest) {Isolation Forest\\[-1pt]\scriptsize anomaly score};
\node[bluebox, below=0.4cm of iforest, minimum width=2.8cm] (markov) {Markov Model\\[-1pt]\scriptsize transition prob};
\node[bluebox, below=0.4cm of markov, minimum width=2.8cm] (defense) {Defence Module\\[-1pt]\scriptsize block / pass};

% Blue team container
\begin{scope}[on background layer]
\node[draw=blue!60, fill=blue!3, rounded corners, fit=(feat)(iforest)(markov)(defense),
      inner sep=0.35cm, label={[font=\small\bfseries, text=blue!70!black]above:Blue Team (Detector + Defence)}] (bluegroup) {};
\end{scope}

% ---- Outputs (far right) ----
\node[outbox, right=1.8cm of feat, yshift=-0.2cm] (metrics) {Metrics JSON\\[-1pt]\scriptsize P, R, F1, latency};
\node[outbox, below=0.3cm of metrics] (plot) {3-Panel Graph\\[-1pt]\scriptsize accuracy, score, blocks};
\node[outbox, below=0.3cm of plot] (logs) {CSV Logs\\[-1pt]\scriptsize defence, alerts};

% Output container
\begin{scope}[on background layer]
\node[draw=yellow!70!black, fill=yellow!4, rounded corners, fit=(metrics)(plot)(logs),
      inner sep=0.3cm, label={[font=\small\bfseries, text=yellow!50!black]above:Outputs}] (outgroup) {};
\end{scope}

% ---- Arrows ----

% Orchestrator controls
\draw[arr] (orch.south) -- ++(0,-0.4) -| node[label, pos=0.25, above]{starts / stops} (sim.north);
\draw[arr] (orch.south) -- ++(0,-0.4) -| (redgroup.north);
\draw[arr] (orch.south) -- ++(0,-0.4) -| (bluegroup.north);
\draw[darr] (orch.south) -- ++(0,-0.4) -| node[label, pos=0.25, above]{collects} (outgroup.north);

% Simulator to UDP
\draw[arr, green!50!black] (sim.south) |- node[label, pos=0.7, above]{\scriptsize normal traffic} (udp.west |- sim.south);
% Adjust: sim sends to udp
\draw[arr, green!50!black] (sim.east) -- ++(0.5,0) |- (udp.north west);

% Attacks to UDP
\draw[arr, red!60!black] (redgroup.east) -- node[label, above]{\scriptsize malicious traffic} (udp.west);

% UDP to Blue Team
\draw[arr] (udp.east) -- node[label, above]{\scriptsize mixed stream} (bluegroup.west |- udp.east);

% Blue team internal flow
\draw[arr] (feat.south) -- (iforest.north);
\draw[arr] (iforest.south) -- (markov.north);
\draw[arr] (markov.south) -- (defense.north);

% Blue team to outputs
\draw[arr] (defense.east) -| ++(0.6, 0) |- (outgroup.west |- plot.west);

% ---- Step labels ----
\node[font=\footnotesize\bfseries, fill=white, draw=gray, circle, inner sep=1.5pt] at ([xshift=-0.5cm, yshift=0.3cm]sim.north west) {1};
\node[font=\footnotesize\bfseries, fill=white, draw=gray, circle, inner sep=1.5pt] at ([xshift=-0.5cm, yshift=0.3cm]redgroup.north west) {2};
\node[font=\footnotesize\bfseries, fill=white, draw=gray, circle, inner sep=1.5pt] at ([xshift=-0.3cm]udp.north west) {3};
\node[font=\footnotesize\bfseries, fill=white, draw=gray, circle, inner sep=1.5pt] at ([xshift=-0.5cm, yshift=0.3cm]bluegroup.north west) {4};
\node[font=\footnotesize\bfseries, fill=white, draw=gray, circle, inner sep=1.5pt] at ([xshift=-0.5cm, yshift=0.3cm]outgroup.north west) {5};

\end{tikzpicture}%
}

\caption{Framework architecture overview.
\textbf{Step~1:} The orchestrator launches the platform simulator, which emits normal MAVLink traffic.
\textbf{Step~2:} Red-team attack modules inject malicious traffic into the same UDP channel.
\textbf{Step~3:} Normal and malicious packets are mixed on UDP port~14550.
\textbf{Step~4:} The blue-team pipeline extracts 11 features per message, scores each message with the Isolation Forest and Markov model, and the defence module decides to block or pass.
\textbf{Step~5:} The orchestrator collects all results into standardised metrics, graphs, and logs.
The orchestrator (top) coordinates all components without knowing their internal implementation.}
\label{fig:framework}
\end{figure*}


\subsection{Abstract Interfaces}

To ensure extensibility, all components implement abstract base classes defined in \texttt{framework/base.py}:

\begin{itemize}
    \item \texttt{BasePlatform}: requires \texttt{start()}, \texttt{stop()}, and a \texttt{name} property.
    \item \texttt{BaseAttack}: requires \texttt{execute(target, duration, rate, out\_dir)}, \texttt{stop()}, \texttt{name}, and \texttt{description}.
    \item \texttt{BaseDetector}: requires \texttt{fit(feature\_matrix)}, \texttt{score(feature\_vector)}, and \texttt{is\_trained}.
    \item \texttt{BaseDefense}: requires \texttt{evaluate(timestamp, msg\_id, src, score)} returning a block decision.
\end{itemize}

Adding a new attack or defence requires implementing only the relevant interface and adding it to the registry (\texttt{attacks/\_\_init\_\_.py}). The orchestrator discovers and runs all registered modules without modification to its own code.

\subsection{Red-Team Attack Catalogue}

Table~\ref{tab:attacks} summarises the six implemented attacks.
Each attack implements the \texttt{BaseAttack} interface and is registered in the central \texttt{ATTACK\_REGISTRY}. The orchestrator discovers and executes them via the \texttt{execute()} method without hard-coding script paths.
Attacks use distinct source system IDs (240--255) so that the detector can observe traffic from multiple origins.

\begin{table}[t]
\centering
\caption{Red-team attack catalogue with default parameters.}
\label{tab:attacks}
\renewcommand{\arraystretch}{1.15}
\begin{tabular}{lll}
\toprule
\textbf{Attack} & \textbf{Message type} & \textbf{Mechanism} \\
\midrule
Heartbeat flood & HEARTBEAT & High-rate flood from rogue GCS \\
Ping flood & PING & Rapid ping requests \\
Param-request flood & PARAM\_REQUEST & Excessive parameter queries \\
MITM identity spoof & HEARTBEAT & Forged system ID~1 messages \\
Replay pattern & PING, PARAM, STATUS & Fixed stale sequence replay \\
Command injection & COMMAND\_LONG & Unauthorised flight commands \\
\bottomrule
\end{tabular}
\end{table}

\subsection{Feature Engineering} \label{sec:features}

The detector extracts an 11-dimensional feature vector for every incoming message using a sliding time window (default 2~seconds).
Features are computed in real time by the \texttt{FeatureEngine} class:

\begin{table}[t]
\centering
\caption{Eleven features extracted per MAVLink message.}
\label{tab:features}
\renewcommand{\arraystretch}{1.15}
\begin{tabular}{clp{4.2cm}}
\toprule
\textbf{\#} & \textbf{Feature} & \textbf{Description} \\
\midrule
1 & \texttt{msg\_rate} & Messages per second in window \\
2 & \texttt{type\_entropy} & Shannon entropy of message-type distribution \\
3 & \texttt{src\_system\_count} & Distinct source system IDs \\
4 & \texttt{src\_msg\_rate} & Rate from current source only \\
5 & \texttt{markov\_prob} & Transition probability (prev $\to$ current type) \\
6 & \texttt{inter\_arrival} & Time since previous message \\
7 & \texttt{ratio\_heartbeat} & Fraction of HEARTBEAT messages \\
8 & \texttt{ratio\_ping} & Fraction of PING messages \\
9 & \texttt{ratio\_param\_req} & Fraction of PARAM\_REQUEST messages \\
10 & \texttt{ratio\_cmd\_long} & Fraction of COMMAND\_LONG messages \\
11 & \texttt{ratio\_statustext} & Fraction of STATUSTEXT messages \\
\bottomrule
\end{tabular}
\end{table}

The combination of rate-based features (1, 4, 6), diversity features (2, 3), sequential features (5), and per-type ratio features (7--11) enables the detector to identify both volumetric floods (elevated rate, low entropy) and subtle injection attacks (anomalous type ratios, unexpected transitions).

\subsection{Anomaly Detection: Isolation Forest}

The primary anomaly detector is an Isolation Forest~\cite{liu2008isolation} trained during the baseline phase on feature vectors from normal traffic.
The model assigns each subsequent message a normalised anomaly score $s \in [0, 1]$, where $s = 1$ indicates normal behaviour and $s = 0$ indicates maximal anomaly.
Score normalisation uses the decision-function range observed during training:

\begin{equation}
    s = \frac{d(x) - d_{\min}}{d_{\max} - d_{\min}}
\end{equation}

\noindent where $d(x)$ is the raw Isolation Forest decision function value and $d_{\min}$, $d_{\max}$ are the minimum and maximum values observed on the training set.

\subsection{Markov Transition Model}

A first-order Markov chain captures the expected message-type transition probabilities from baseline traffic.
During monitoring, the transition probability $P(m_t \mid m_{t-1})$ is computed for each incoming message type $m_t$ given the previous type $m_{t-1}$:

\begin{equation}
    P(m_t \mid m_{t-1}) = \frac{C(m_{t-1}, m_t)}{\sum_{j} C(m_{t-1}, j)}
\end{equation}

\noindent where $C(i, j)$ is the count of observed transitions from type $i$ to type $j$ during baseline.
A low transition probability indicates an unexpected message sequence, which is characteristic of injected attack traffic.

\subsection{Defence Strategies}

Two active defence mechanisms are implemented:

\textbf{Markov Defence} (\texttt{MarkovDefense}):
Blocks any message whose transition probability $P(m_t \mid m_{t-1})$ falls below a configurable threshold $\tau_M$ (default 0.05).
This is a lightweight, pattern-based defence that uses the learned message ordering.

\textbf{Adaptive Defence} (\texttt{AdaptiveDefense}):
Blocks any message whose Isolation Forest anomaly score falls below threshold $\tau_A$ (default 0.5).
Because the Isolation Forest normalises its decision function to $[0, 1]$ using the training-set range (Equation~1), a threshold of 0.5 naturally separates the lower half of the score distribution from the upper half, providing a principled default.
This defence uses the full 11-feature vector and adapts to attacks that evade simple sequential patterns.

\textbf{Rate-Threshold Baseline} (\texttt{RateDefense}):
A deliberately simple defence included as a lower-bound comparison.
It learns the median message rate $\bar{r}$ from the baseline phase and blocks any message arriving when the instantaneous rate exceeds $\bar{r} \times k$ for a configurable multiplier $k$ (default 2).
Because it relies on a single feature (message rate), it cannot distinguish between different attack types or detect low-rate injection attacks.

All three defences log every decision (block/pass) with timestamp, message ID, source system, and the score that triggered the decision.

\subsection{Orchestrator Operation}

The orchestrator (\texttt{run\_demo.py}) manages the complete experimental lifecycle.
Algorithm~\ref{alg:orchestrator} presents the orchestration logic.

\begin{algorithm}[t]
\caption{Orchestrator: evaluation scenario execution}\label{alg:orchestrator}
\DontPrintSemicolon
\KwIn{Configuration $\mathcal{C}$: UDP address, sim rate, baseline duration $T_b$, attack schedule $\mathcal{A}$, defence mode $\delta$, thresholds}
\KwOut{Metrics $\mathcal{M}$, plots, CSV logs}
\BlankLine
Launch platform simulator at configured rate\;
Initialise detector (Isolation Forest + Markov model)\;
Start detector in background thread\;
\BlankLine
\tcp{Phase 1: Baseline training}
Wait $T_b$ seconds (detector collects normal traffic)\;
Detector trains Isolation Forest on baseline features\;
Detector builds Markov transition matrix\;
\BlankLine
\tcp{Phase 2: Attack injection}
\ForEach{episode $e_i \in \mathcal{A}$}{
    \If{randomised mode}{
        Wait random inter-episode gap $g_i$\;
    }
    \ForEach{segment $(r_j, d_j) \in e_i.\text{segments}$}{
        Execute attack plugin with rate $r_j$ for duration $d_j$\;
        Detector scores each message, defence blocks if below threshold\;
        Log attack start/end timestamps\;
    }
}
\BlankLine
\tcp{Phase 3: Result collection}
Stop simulator\;
Compute classification metrics (precision, recall, F1)\;
Compute mission metrics (detection latency, block rate, accuracy recovery)\;
Generate three-panel diagnostic plot\;
Write JSON results and CSV logs\;
\Return $\mathcal{M}$\;
\end{algorithm}

\subsection{Randomised Attack Scheduling}

In randomised mode (\texttt{--randomize-attacks}), the orchestrator transforms the static attack schedule into a volatile, unpredictable pattern:

\begin{enumerate}
    \item The attack order is shuffled uniformly at random.
    \item Each episode duration is sampled from $[\,0.4 \cdot T_a,\; 1.6 \cdot T_a\,]$ where $T_a$ is the base attack duration.
    \item Each episode is split into 5--13 segments with independently sampled rates from $[\,0.35 \cdot R,\; 1.65 \cdot R\,]$ where $R$ is the base attack rate.
    \item Random micro-pauses (0.04--0.75\,s) are inserted between segments.
    \item Random inter-episode gaps (0.35--5.5\,s) separate consecutive attacks.
\end{enumerate}

This scheduling models a realistic adversary who varies attack intensity and timing to evade pattern-based defences.
The moving-average window for accuracy computation is automatically capped at 160 messages to ensure the plot reflects rapid traffic changes rather than smoothing them out.

\subsection{Metrics and Evaluation Criteria}

The framework computes the following standardised metrics for each experimental run:

\begin{itemize}
    \item \textbf{Detection rate} (recall): fraction of attack-period messages correctly flagged.
    \item \textbf{False-positive rate}: fraction of baseline alerts that occur outside attack windows.
    \item \textbf{Detection latency}: time from first attack onset to first alert.
    \item \textbf{Block rate}: fraction of messages actively blocked by the defence.
    \item \textbf{Precision}: proportion of blocked/alerted messages that were truly attacks.
    \item \textbf{F1 score}: harmonic mean of precision and recall.
    \item \textbf{Mission impact ratio}: fraction of the monitoring period with accuracy below 0.9.
    \item \textbf{Accuracy recovery}: difference between post-attack maximum accuracy and attack-period minimum.
\end{itemize}

All metrics are computed by the \texttt{ScenarioResult} class and exported as JSON for automated comparison across configurations.
